\section{Worked traces}
\label{app:tables}

The state set and the transition rules are given in Tables~\ref{tab:states} to~\ref{tab:reg}. Here we show two
short traces, to make the rules concrete.

\subsection{The locomotive on a straight track}

Table~\ref{tab:loco-trace} traces the locomotive along five track cells, written left to right as the head
moves. Each row is one step. The head $H$ advances by one cell, the trailing cell $A$ follows, and the cells
behind clear to $C$, exactly by the motion rules of Table~\ref{tab:motion}.

\begin{table}[ht]
\centering
\begin{tabular}{@{}cccccc@{}}
\toprule
step & cell 1 & cell 2 & cell 3 & cell 4 & cell 5 \\
\midrule
0 & $A$ & $H$ & $C$ & $C$ & $C$ \\
1 & $C$ & $A$ & $H$ & $C$ & $C$ \\
2 & $C$ & $C$ & $A$ & $H$ & $C$ \\
3 & $C$ & $C$ & $C$ & $A$ & $H$ \\
\bottomrule
\end{tabular}
\caption{The locomotive advancing along a straight track. The head and its trailing cell move forward together,
one cell per step, and never backward.}
\label{tab:loco-trace}
\end{table}

\subsection{The self-extending binary counter}

Table~\ref{tab:counter-trace} traces the self-extending counter of Section~\ref{sec:strong} through its first
eight increments. Each row gives the bits after the increment, lowest bit on the right, with a dot for a bit the
counter has not yet built. The counter starts as a single bit and builds a new one each time the count first
reaches a power of two.

\begin{table}[ht]
\centering
\begin{tabular}{@{}ccccl@{}}
\toprule
count & bit 3 & bit 2 & bit 1 & bit 0 \\
\midrule
0 & $\cdot$ & $\cdot$ & $\cdot$ & $0$ \\
1 & $\cdot$ & $\cdot$ & $\cdot$ & $1$ \\
2 & $\cdot$ & $\cdot$ & $1$ & $0$ \quad (built bit 1) \\
3 & $\cdot$ & $\cdot$ & $1$ & $1$ \\
4 & $\cdot$ & $1$ & $0$ & $0$ \quad (built bit 2) \\
5 & $\cdot$ & $1$ & $0$ & $1$ \\
6 & $\cdot$ & $1$ & $1$ & $0$ \\
7 & $\cdot$ & $1$ & $1$ & $1$ \\
8 & $1$ & $0$ & $0$ & $0$ \quad (built bit 3) \\
\bottomrule
\end{tabular}
\caption{The self-extending counter. A dot marks a bit not yet built. The counter begins with bit 0 alone and
builds bit $n$ exactly when the count first reaches $2^n$, by the track-building rule of Table~\ref{tab:build}.
The memory grows on demand, from a finite seed.}
\label{tab:counter-trace}
\end{table}
